\documentclass[11pt]{article}
\usepackage[margin=1in]{geometry}
\usepackage{booktabs}
\usepackage{hyperref}
\usepackage{url}

\title{Memo 13: MP-only Numerical Limiter Update for \texttt{PLAST\_FLOC=F}\\
       Baseline Stability in \texttt{mod\_plast.F}}
\author{FVCOM-MP WaterPACT Investigation}
\date{May 25, 2026}

\begin{document}
\sloppy
\maketitle

\section*{Purpose}

This memo records the model-development change made after the MP-only
baseline runs with \texttt{PLAST\_FLOC=F} showed numerical divergence.  The
main purpose of the change is to make the non-floc plastic pathway mirror the
same water-column upper-bound value limiter already used by the split
\texttt{PLAST\_FLOC=T} pathway.

The project-level reference consulted for this memo was:
\begin{quote}
\path{C:\Users\huan111\OneDrive - PNNL\Desktop\WaterPACT_Local\4_Floc_MP_interaction\Memory\v003_mp_mass_divergence_diagnosis.tex}
\end{quote}

\section*{Background}

The divergent cases were MP-only baseline tests using
\texttt{PLAST\_FLOC=F}.  The reference memo showed that the river MP inputs
were small and finite, while the model plastic mass grew to unrealistically
large values during the run.  The growth was therefore not consistent with a
forcing-magnitude problem.  Instead, the diagnostic checkpoints suggested that
the non-floc transport pathway could create very large positive or negative
combined plastic concentrations before the later lower-bound clipping stage.

The important numerical asymmetry was:
\begin{itemize}
  \item With \texttt{PLAST\_FLOC=T}, the split plastic fields
  \texttt{cnew\_a} and \texttt{cnew\_d} were already capped at 100 before the
  \texttt{CP7\_after\_upper\_clamp} mass checkpoint.
  \item With \texttt{PLAST\_FLOC=F}, the combined non-floc field
  \texttt{cnew} passed through advection and vertical diffusion without the
  corresponding upper-bound cap.
\end{itemize}

That difference made the nominal MP-only baseline a different numerical
pathway, not simply a floc-disabled version of the coupled pathway.  In
particular, once very large negative values were generated, the later
non-negative clamp could zero them while leaving compensating positive
excursions, producing an artificial mass increase.

\section*{Code Change Recorded}

The current development version of \texttt{mod\_plast.F} adds an upper-bound
limiter to the combined non-floc concentration:
\begin{verbatim}
where (plast(iplast)%cnew > 100.0_sp) plast(iplast)%cnew = 100.0_sp
\end{verbatim}

This limiter is applied after the non-floc advection and vertical diffusion
updates and before \texttt{CP7\_after\_upper\_clamp}.  The same treatment was
added in both relevant non-floc branches:
\begin{itemize}
  \item the sediment-enabled branch used by the main FVCOM-MP sediment build;
  \item the non-sediment fallback branch guarded by the preprocessor
  alternative.
\end{itemize}

The mass-check labels were also clarified in the non-floc branch.  For
\texttt{PLAST\_FLOC=F}, checkpoints \texttt{CP3p5} through \texttt{CP6}
represent the same post-transport, pre-cap combined concentration state.
\texttt{CP7\_after\_upper\_clamp} now records the state after the new
upper-bound limiter.

\section*{Before and After}

\begin{center}
\begin{tabular}{p{0.24\linewidth}p{0.34\linewidth}p{0.34\linewidth}}
\toprule
Aspect & Before & After \\
\midrule
\texttt{PLAST\_FLOC=T} split path &
\texttt{cnew\_a} and \texttt{cnew\_d} capped at 100 before
\texttt{CP7}. &
Unchanged. \\
\addlinespace
\texttt{PLAST\_FLOC=F} combined path &
\texttt{cnew} was advected and vertically diffused, but no equivalent
upper cap was applied before \texttt{CP7}. &
\texttt{cnew} is capped at 100 before
\texttt{CP7\_after\_upper\_clamp}. \\
\addlinespace
Lower-bound treatment &
The later non-negative clamp was already applied after open-boundary and
point-source handling. &
Unchanged.  The new change only adds the missing upper-bound limiter. \\
\addlinespace
Physical parameterization &
No change to aggregation, disaggregation, settling, density, or floc
coupling physics. &
Still no change to those physics.  This is a numerical value limiter. \\
\bottomrule
\end{tabular}
\end{center}

\section*{Interpretation}

This update makes the MP-only baseline numerically more comparable to the
floc-enabled configuration.  It does not turn on floc coupling, aggregation,
disaggregation, equivalent grain size interaction, or sediment-floc physics
when \texttt{PLAST\_FLOC=F}.  It only prevents the combined non-floc plastic
field from exceeding the same upper concentration bound that already protected
the split \texttt{PLAST\_FLOC=T} fields.

The change should reduce the risk that an MP-only baseline diverges simply
because it lacks a limiter present in the coupled run.  Therefore, if the
baseline becomes stable after this update, the result should be interpreted as
a correction to the numerical pathway rather than a change in the physical
meaning of the baseline case.

\section*{Recommended Verification Runs}

For the next short P0/P1 verification, use the same forcing and model setup
that previously diverged, but retain the plastic mass-check output.  The key
diagnostic comparison is:
\begin{itemize}
  \item \texttt{CP3\_after\_update\_bottom}: mass before water-column
  advection and diffusion;
  \item \texttt{CP3p5\_after\_adv\_a} through \texttt{CP6\_after\_kinetics}:
  the post-transport, pre-cap combined non-floc state;
  \item \texttt{CP7\_after\_upper\_clamp}: the first checkpoint after the
  new cap;
  \item \texttt{CP10\_after\_neg\_clamp}: the state after the existing
  non-negative lower-bound clamp.
\end{itemize}

The expected healthy behavior is that \texttt{CP7} no longer permits extreme
positive concentrations, and \texttt{CP10} no longer shows explosive mass
growth caused by clipping very large negative values.  If divergence persists,
the next isolation tests should be a transport-only run with point sources
disabled, followed by a focused search for the first cell and time step where
\texttt{cnew} becomes non-finite or very large in magnitude.

\section*{Summary}

The development repository now includes a missing upper-bound limiter for the
combined MP concentration used when \texttt{PLAST\_FLOC=F}.  This brings the
MP-only pathway into closer numerical alignment with the existing
\texttt{PLAST\_FLOC=T} split-field pathway and should make the baseline tests a
fairer comparison for long-run MP-floc experiments.

\end{document}
